\section*{Chapitre \ref{ch:in:quant-developer} --- Développeur quant}

\begin{solution}{exo:in:quant-developer:1}
122 sur 169, 72.2~\%, comme analystes quantitatifs financiers~; 25 sur 169, 14.8~\%, comme développeurs de logiciels~; les 13.0~\% restants sous
des codes comptant chacun moins de dix dépôts.
\end{solution}

\begin{solution}{exo:in:quant-developer:2}
(\texttt{QUANTITATIVE DEVELOPER}, \texttt{senior})~: la parenthèse et ce qui la suit sont supprimées, «~Sr~» est enregistré comme
rang de séniorité et «~Quant Dev~» est développé.
\end{solution}

\begin{solution}{exo:in:quant-developer:3}
$163\,290/136\,620=1.195$~: environ un cinquième de plus.
\end{solution}

\begin{solution}{exo:in:quant-developer:4}
La médiane globale de chaque code mélange ses niveaux. La moitié des dépôts de développeurs logiciels sont au niveau IV, contre un sur
sept pour les analystes. La médiane globale du code logiciel est donc tirée vers le haut par ses dépôts de niveau supérieur, alors que, niveau
par niveau, les dépôts des analystes portent l'offre la plus élevée.
\end{solution}

\begin{solution}{exo:in:quant-developer:5}
Avec la répartition des analystes, les médianes par niveau des développeurs logiciels sont en moyenne inférieures de 7\,789~\$ à celles des analystes~: le signe de
l'écart global est une question de composition.
\end{solution}

\begin{solution}{exo:in:quant-developer:6}
Le poste d'\omterm{def:in:quant-developer:research}{ingénieur de recherche}, qui travaille au sein d'un groupe de recherche et fait souvent passer des signaux en production.
Elle devrait demander si les ingénieurs y traitent leurs propres problèmes de recherche, combien sont passés à la recherche,
et qui possède les signaux qu'elle devrait mettre en œuvre.
\end{solution}

\begin{solution}{exo:in:quant-developer:7}
Cinq sur dix~: pour l'exercice fiscal 2021, l'écart global et les niveaux III et IV~; pour l'exercice fiscal 2025, les niveaux III et IV. Le même écart de l'exercice fiscal 2025
L'écart global ($-\$372$ à 5\,688~\$) et chaque écart des niveaux I et II incluent zéro.
\end{solution}

\begin{solution}{exo:in:quant-developer:8}
Trois défauts. L'intervalle de l'écart de 5\,000~\$ contient zéro. Il compare des populations aux répartitions de niveaux différentes~; à l'intérieur de chaque
niveau, la médiane des analystes est au moins aussi élevée. Enfin, le code décrit le poste pour lequel l'employeur dépose une déclaration~; il ne fixe pas
l'offre, qui dépend de l'employeur et du niveau.
\end{solution}

\begin{solution}{pb:in:quant-developer:1}
\begin{enumerate}
\item Un ingénieur qui transforme les prototypes de recherche en code de production et répond de sa performance et de sa justesse~;
l'\omterm{def:in:quant-developer:research}{ingénieur de recherche}, le \omterm{def:in:quant-developer:library}{développeur de bibliothèque} et le \omterm{def:in:quant-developer:trading}{développeur de systèmes de trading}.
\item Infrastructure de recherche et signaux de production~; versions de bibliothèques~; le système de trading.
\item Reproductibilité et temps des chercheurs~; stabilité, rapidité et couverture~; latence, disponibilité et comportement correct. Un résultat
que personne ne peut rejouer~; les chiffres de chaque desk qui bougent~; une perte en quelques secondes.
\item Le \omterm{def:in:quant-developer:trading}{développeur de systèmes de trading}.
\item Dans un groupe de recherche~; dans une équipe centrale d'analyse (ou la technologie de trading d'un teneur de marché d'options)~;
dans la technologie de trading.
\item Les développeurs de logiciels et les analystes quantitatifs financiers~; les premiers correspondent au \omterm{def:in:quant-developer:trading}{développeur de systèmes de trading}, les seconds
à l'\omterm{def:in:quant-developer:research}{ingénieur de recherche} et au \omterm{def:in:quant-developer:library}{développeur de bibliothèque}.
\item Le code des analystes se situe dans la tranche de préparation la plus élevée (études de troisième cycle), celui des développeurs de logiciels dans la tranche
inférieure (premier cycle universitaire).
\item Les codes 2010 15-1132 et 15-1133 ont été fusionnés dans le 15-1252.
\item 72.2~\% 13-2099.01, 14.8~\% 15-1252, 13.0~\% autres codes.
\item Presque tous les dépôts proviennent d'un ou deux employeurs par catégorie, et les cellules issues de moins de trois employeurs
sont supprimées.
\item Rééchantillonnages bootstrap indépendants de chaque population~; la différence et le rapport des médianes rééchantillonnées donnent des intervalles à
95\up{e}~centile.
\item 5\,000~\$, de $-\$372$ à 5\,688~\$.
\item $-\$11\,000$ ($-\$19\,038$ à $-\$6\,600$) et $-\$27\,092$ ($-\$37\,500$ à $-\$19\,900$).
\item Les dépôts du code logiciel sont concentrés au niveau IV.
\item Il devient $-\$7\,789$.
\item 72.2~\% analystes, 14.8~\% développeurs de logiciels, 13.0~\% autres~; logiciels moins analystes 5\,000~\$ au total ($-\$372$ à
5\,688~\$), $-\$27\,092$ au niveau IV.
\item Non~: l'employeur et le niveau fixent l'offre~; le code ne fait que la décrire.
\item Une médiane de 235\,000~\$ sur 12 dépôts issus de 3 employeurs~; les dépôts des autres catégories proviennent d'un ou deux employeurs.
\item Ce que livre le poste~: des outils de recherche, une bibliothèque ou un système de trading, et l'astreinte qui va avec.
\item Les trois offres sont trois postes différents (outils de recherche, bibliothèque de pricing, système de trading), quels que soient leur intitulé
et leur code. Choisir selon le produit et le rythme auquel on veut travailler, et comparer la rémunération au sein de la catégorie d'employeur et du niveau.
\end{enumerate}
\end{solution}

\begin{solution}{iq:in:quant-developer:1}
Déterminer s'il s'agit d'arithmétique (ordre des opérations, sommation en virgule flottante, fonction de bibliothèque différente) ou de données (une
entrée différente ces jours-là)~: consigner les valeurs intermédiaires des deux versions un jour en échec et les comparer pas à pas~; puis
fixer la tolérance avec le chercheur et ajouter le cas aux tests de non-régression.
\iqlookfor{dichotomie sur le calcul, et tolérance convenue et testée.}
\end{solution}

\begin{solution}{iq:in:quant-developer:2}
Versionner le code, les données (instantanés point-in-time ou version des données), la configuration et l'environnement
(versions de dépendances figées, graines aléatoires), les consigner avec le résultat, et rejouer le calcul à partir de cet enregistrement comme test.
\iqlookfor{code, données, configuration et environnement tous versionnés.}
\end{solution}

\begin{solution}{iq:in:quant-developer:3}
Profiler avant toute modification~: allocations et copies à chaque appel, travail répété (courbes reconstruites à chaque appel), défauts de cache (cache miss),
puis la méthode numérique~; regrouper les appels par transactions~; paralléliser sur les cœurs en dernier.
\iqlookfor{mesurer d'abord~; travail répété et mémoire avant l'arithmétique.}
\end{solution}

\begin{solution}{iq:in:quant-developer:4}
Comparer les histogrammes de latence des deux versions sur les mêmes données rejouées~; procéder par dichotomie sur les modifications de la version~; chercher de nouvelles
allocations, verrous, journalisation ou branchements sur un chemin rare, et les changements de placement des threads ou de ramasse-miettes (garbage collection).
\iqlookfor{la latence de queue demande rejeu et dichotomie, pas des moyennes.}
\end{solution}

\begin{solution}{iq:in:quant-developer:5}
Une spécification typée des entrées avec leurs horodatages, les unités et la plage de sortie, des cas de référence avec les sorties attendues
que les deux versions exécutent en CI, et un contrôle en direct qui compare les deux sur des données miroir.
\iqlookfor{cas de référence communs et comparaison en direct.}
\end{solution}

\begin{solution}{iq:in:quant-developer:6}
Valeurs de référence (golden values) pour un portefeuille figé, comparées à la version précédente et à un indice de référence indépendant, avec des tolérances par produit~;
identités (parité, bornes) et grecques par rapport aux différences finies~; tests de performance~; chaque bogue passé en test~; et un
rapport de chaque chiffre modifié à l'intention des desks avant la mise en production.
\iqlookfor{non-régression sur valeurs de référence avec tolérances, références indépendantes et rapport des modifications.}
\end{solution}
